\documentclass[10pt,letterpaper]{article}
\usepackage{cornell-notes}

\title{Clause 05.02.11: Ad Element Correspondence}
\author{}
\date{}
\setDocTitle{Clause 05.02.11: Ad Element Correspondence}
\setDocSubtitle{ISO/IEC/IEEE 42010:2022}
\setDocOwner{ISO 42010 Cornell Notes}
\setDocDate{\today}
\setCornellCollection{ISO/IEC/IEEE 42010:2022 Cornell Notes}
\setCornellUnitType{Clause}
\setCornellUnitNumber{05.02.11}
\setCornellUnitTitle{Ad Element Correspondence}
\hypersetup{pdftitle={Clause 05.02.11: Ad Element Correspondence},pdfsubject={Cornell notes},pdfauthor={}}

\begin{document}
\maketitle
\begin{CornellOverviewBox}{Overview}
{\Large\bfseries\textcolor{CNavy}{Section 5.2.11: AD element correspondence}}\\[3pt]
{\small\textbf{Classification:} Conceptual foundation}\\
{\small\textbf{Learning objective:} Explain the role of ad element correspondence and connect it to related architecture-description concepts.}
\end{CornellOverviewBox}
\vspace{3pt}

\begin{CornellNotesTable}
\CornellNoteRow{Purpose / key concept}{An AD element correspondence identifies an identified or named relation between two or more AD elements.}
\CornellNoteRow{Distinction}{A correspondence is the relation itself; a correspondence method specifies rules or practices governing such relations.}
\CornellNoteRow{Study relationship / visual insight}{Study relationship: a correspondence relates one or more AD elements; a correspondence method governs that relation; for correspondence purposes, an architecture description can itself participate as an AD element in another description.}
\CornellNoteRow{Example / application}{A correspondence between an AD element within a view and the concern that it addresses; between an architecture view and the aspect that it implements; between an AD element and the function that it implements; between an interface on a component and the stack of standards to which the interface conforms; between a data object on a functional flow and the full data structure definition; between a system and the organizational structure that implements it.}
\CornellNoteRow{Interpretive note}{\begin{itemize}[leftmargin=*,nosep,topsep=1pt]
\item Requirements for using correspondences and correspondence methods are specified in 6.9. Additional
\item Correspondences Here are similar to view correspondences in ISO/IEC 10746-2 (RM- ODP) and ISO/IEC 19793.
\end{itemize}}
\CornellNoteRow{Related sections}{6.9}
\CornellNoteRow{Active recall}{\begin{itemize}[leftmargin=*,nosep,topsep=1pt]
\item What is the central role of AD element correspondence?
\item How does AD element correspondence connect to adjacent architecture-description concepts?
\item What closely related concepts must be kept distinct?
\end{itemize}}
\end{CornellNotesTable}


\begin{CornellSummaryBox}{Summary}
An AD element correspondence identifies an identified or named relation between two or more AD elements. A correspondence is the relation itself; a correspondence method specifies rules or practices governing such relations.
\end{CornellSummaryBox}

\end{document}
